\section{The inherited physical interface}
\label{sec:tpc43-interface}

We use only the terminal statements proved in TPC--18, TPC--36,
TPC--41, and TPC--42
\citep{WangTPC18,WangTPC36,WangTPC41,WangTPC42}.  Fix
\(h\in\Z\setminus\{0\}\).  A source row is
\begin{equation}
 \alpha=(\ell_\alpha,d_\alpha),
 \qquad
 m_\alpha=\ell_\alpha d_\alpha\asymp Q,
 \qquad
 \ell_\alpha\asymp L_0,
 \qquad
 d_\alpha\asymp D_0,
 \label{eq:tpc43-row}
\end{equation}
where \(\ell_\alpha\) is prime, \(d_\alpha\) is squarefree,
\((d_\alpha,h)=1\), and \(L_0D_0=Q\).  Source separation makes
\(m_\alpha\) determine \(\alpha\).  The orbit variable belongs to one
of the inherited dyadic blocks
\begin{equation}
 a_JJ\le j\le b_JJ,
 \qquad QJ\asymp X,
 \label{eq:tpc43-orbit-block}
\end{equation}
for fixed positive \(a_J<b_J\).  At the endpoint,
\begin{align}
 Q&=X^{267/400+o(1)},
 &J&=X^{133/400+o(1)},
 \label{eq:tpc43-QJ-endpoint}\\
 D_0&=X^{10049/52500+o(1)}=o(J).
 \label{eq:tpc43-D-endpoint}
\end{align}
In particular, for all sufficiently large \(X\),
\begin{equation}
 \min\cJ>|h|\max_\alpha d_\alpha.
 \label{eq:tpc43-rigidity-size}
\end{equation}

Let \(M\) be the inherited triple--prime modulus and let
\(\cI_B=\{-L_B,\ldots,L_B\}\), where
\begin{equation}
 K_B:=|\cI_B|=X^{1/400+o(1)},
 \qquad
 K_B=X^{o(1)}\frac{Q}{J^2}.
 \label{eq:tpc43-alias-ledger}
\end{equation}
For \(k\in\cI_B\), put
\begin{equation}
 c_k=h+kM,
 \qquad
 n_{\alpha,j,k}=m_\alpha j+c_k.
 \label{eq:tpc43-terminal-n}
\end{equation}
Because \(h\) is fixed and nonzero while \(M>|h|\) for all sufficiently
large \(X\), every \(c_k\) occurring here is nonzero: this is immediate
for \(k=0\), and \(|kM|>|h|\) for \(k\ne0\).
All terminal coefficients are interpreted as zero outside the declared
positive support.  Uniformly on that support,
\begin{equation}
 |c_k|+d_\alpha n_{\alpha,j,k}\le X^{C_0}
 \label{eq:tpc43-polynomial-support}
\end{equation}
for one fixed \(C_0\).

The source coefficient factors as
\begin{equation}
 \gamma_\alpha^{(i)}
 =\mu(d_\alpha)\gamma_\alpha^{\circ,(i)},
 \qquad i=1,2,
 \label{eq:tpc43-strip-source-mobius}
\end{equation}
where \(\gamma^{\circ,(i)}\) contains the source logarithm, smooth
weight, and fixed residue data, but no M\"obius sign.  After also stripping
the target M\"obius sign, define deterministic coefficients
\begin{equation}
 a^L_{\alpha,\gamma,j,k}
 :=-\gamma_\alpha^{\circ,(1)}
   \mathfrak A_{\alpha,\gamma}^{\rm act}(j)
   \log(n_{\alpha,j,k})\Phi_0(n_{\alpha,j,k}),
 \label{eq:tpc43-deterministic-left-atom}
\end{equation}
again set to zero outside the terminal support.  Then the physical left
column is exactly
\begin{equation}
 [\mathscr T_k^L]_{\gamma,j}
 =\sum_\alpha
 a^L_{\alpha,\gamma,j,k}
 \mu(d_\alpha)\mu(n_{\alpha,j,k}).
 \label{eq:tpc43-physical-left}
\end{equation}
The right column has coefficients
\(a^R_{\alpha,\gamma,j,k}\) obtained by exchanging the two row systems.
Conditional on the inherited row geometry and terminal support, no
additional structural assumption is made on the deterministic complex
values \(a^{L/R}\); in particular, they are independent of the
prime--sign environment.

Define the atomic diagonals
\begin{align}
 \cD_{k,L}
 &:=\sum_{\gamma,j,\alpha}
 |a^L_{\alpha,\gamma,j,k}|^2
 \mu^2(n_{\alpha,j,k}),
 \label{eq:tpc43-left-diagonal}\\
 \cD_0&:=\cD_{0,L}+\cD_{0,R},
 &
 \cD_{\rm tr}
 &:=\sum_{k\in\cI_B}(\cD_{k,L}+\cD_{k,R}).
 \label{eq:tpc43-total-diagonals}
\end{align}
TPC--41 and TPC--42 give
\begin{align}
 \cD_0&\ll X^{o(1)}Q^2J,
 \label{eq:tpc43-inherited-zero-diagonal}\\
 \cD_{\rm tr}&\ll X^{o(1)}Q^2JK_B.
 \label{eq:tpc43-inherited-trace-diagonal}
\end{align}
The physical equality target is
\begin{equation}
 \|\mathscr T_0^L\|^2+\|\mathscr T_0^R\|^2
 \ll X^{o(1)}K_BQ^2J,
 \label{eq:tpc43-physical-zero-target}
\end{equation}
while the stronger trace target is
\begin{equation}
 \sum_{k\in\cI_B}
 \bigl(\|\mathscr T_k^L\|^2+\|\mathscr T_k^R\|^2\bigr)
 \ll X^{o(1)}Q^2JK_B.
 \label{eq:tpc43-physical-trace-target}
\end{equation}
These two estimates are not assumed.
